\subsection*{身份应当出现在第一次发送之前}

服务接到两个相同报文时，能看到它们的账户、商品和参数，也能看到到达时间。用户在界面上做过什么，通常只存在于调用方的状态里。Alice 点击购买后，客户端创建一项本地操作，即一条记录本次确认内容、操作编号与当前进度的客户端数据；网络超时引起的再次发送仍然引用这项操作。Alice 在购买完成后再次点击购买，客户端创建另一项操作。这个区分在 HTTP 请求发出以前已经存在。

服务为什么无法从报文内容独立推断用户意图，可以把两条历史并排展开。历史 A 中，Alice 确认一次购买，第一次响应丢失，客户端重发；历史 B 中，Alice 确认两次购买，两次响应都正常。假定客户端每次发送的请求都只有账户、商品和数量，服务在两条历史中可以收到相同的两份业务输入，到达时间间隔也可以恰好相同。服务若一概合并，历史 B 少买一份；服务若一概执行，历史 A 多买一份。

\begin{figure}[htbp]
\centering
\begin{tikzpicture}[x=1cm,y=1cm]
\node[booknote,text width=3.2cm] at (1.8,.7) {调用方发生的动作};
\node[booknote,text width=4cm] at (6.2,.7) {服务收到的输入};
\node[booknote,text width=3.2cm] at (11.1,.7) {期望的业务结果};
\node[booknode,text width=3cm] (a) at (1.8,-.4) {A：确认一次\\响应丢失后重发};
\node[booknode,text width=4cm] (ar) at (6.2,-.4) {账户 1、paper-001、1\\账户 1、paper-001、1};
\node[booknode,text width=3cm] (ae) at (11.1,-.4) {保留一笔购买\\余额 7000};
\node[booknode,text width=3cm] (b) at (1.8,-2) {B：明确确认两次\\两次各发一份请求};
\node[booknode,text width=4cm] (br) at (6.2,-2) {账户 1、paper-001、1\\账户 1、paper-001、1};
\node[bookresult,text width=3cm] (be) at (11.1,-2) {形成两笔购买\\余额 4000};
\draw[bookarrow] (a)--(ar);
\draw[bookarrow] (ar)--(ae);
\draw[bookarrow] (b)--(br);
\draw[bookalert] (br)--(be);
\draw[draw=Accent,dashed] (4,-3.05) rectangle (8.4,.2);
\node[booknote,text width=11.5cm] at (6.2,-3.5)
 {第二次到达前，服务两边都已保存 P1，余额均为 7000；相同输入需要不同处理};
\end{tikzpicture}
\caption{缺少操作关系时，两条需要不同结果的历史可以给服务相同的观察。}
\label{fig:ch03-indistinguishable}
\end{figure}


这个例子还可以把判断时点收窄到第二份请求到达的一刻。两条历史中，数据库都已经有 P1、O1、D1，余额都是 7000；服务都已经发送过第一份成功响应。响应是否到达 Alice，以及 Alice 是否再次明确确认，尚未成为服务可读的业务记录。第二份输入相同，服务保存的状态也相同，要求的处理却不同：A 应返回 P1，B 应创建 P2。

把服务用于判断的信息记为 \(O\)，包括这份输入、此前保存的记录和它能够读取的通信信息。上述两条历史可以满足 \(O_A=O_B\)。任何只依赖 \(O\) 的确定性判断，在两条历史中都会给出同一答案，因此无法同时符合两种购买授权。允许随机选择也无法保证每条历史都正确，只会把某些错误变为有概率发生。这是一项信息不足的结论；它的前提是用户是否再次确认没有通过其他可信通道传给服务。若产品另外保存了购买意图，这份记录本身就增加了可用信息。

调用方提供操作关系以后，历史 A 发送 K1、K1，历史 B 发送 K1、K2，服务才获得区分两条历史的输入。我们把 K1 称为幂等键：调用方为一次业务操作选定、在它的多次发送中保持不变的标识。服务以它查找先前的处理记录。第一次发送使用 K1，等待超时后的重试继续使用 K1；Alice 确认另买一份时，客户端创建 K2。K1 再次到达，要求继续处理或返回原购买；K2 到达，申请受理一项新购买。

\begin{figure}[htbp]
\centering
\begin{tikzpicture}[x=1cm,y=1cm]
\node[booknode,text width=4cm] (u1) at (3,0) {Alice 确认买一份\\本地操作 K1};
\node[booknode,text width=3cm] (r1) at (1,-1.6) {R1 携带 K1\\首次发送};
\node[booknode,text width=3cm] (r2) at (5,-1.6) {R2 携带 K1\\超时后恢复};
\node[booknode,text width=4cm] (p1) at (3,-3.2) {共同指向 P1\\一次购买的原结果};
\draw[bookarrow] (u1)--(r1);
\draw[bookarrow] (u1)--(r2);
\draw[bookarrow] (r1)--(p1);
\draw[bookarrow] (r2)--(p1);
\node[booknode,text width=3.7cm] (u2) at (10,0) {Alice 明确再买一份\\新本地操作 K2};
\node[booknode,text width=3.7cm] (r3) at (10,-1.6) {R3 携带 K2\\另一项购买};
\node[bookresult,text width=3.7cm] (p2) at (10,-3.2) {指向 P2\\允许再产生一次效果};
\draw[bookarrow] (u2)--(r3);
\draw[bookalert] (r3)--(p2);
\end{tikzpicture}
\caption{请求按业务操作归组。重试沿用 K1，第二次明确购买使用 K2。}
\label{fig:ch03-intent-attempts}
\end{figure}


这里的业务意图需要落实成可操作的界面规则。两次点击可能来自手抖，也可能来自两次明确购买。客户端可以在一次确认动作发生后把按钮切换为处理中，使后续点击回到该操作的状态页；购买完成后，再次购买按钮明确创建新操作。页面上已经改变商品、数量或接受价格时，则应让用户重新确认并创建新身份。服务能检查身份和输入是否符合协议，确认动作与身份的对应仍由调用方建立。

为了检查这项对应，可以把每次请求尝试记为 \(r\)，把它所属的已确认操作记为 \(I(r)\)，把随请求携带的完整操作身份记为 \(S(r)\)。这里的 \(I\) 是描述购买授权的记号，服务未必能够直接读取它；\(S\) 是协议实际传递并检查的数据。身份规则需要同时满足：
\[
 I(r)=I(r')\ \Longrightarrow\ S(r)=S(r'),\qquad
 I(r)\ne I(r')\ \Longrightarrow\ S(r)\ne S(r').
\]
第一项要求把同一次购买的尝试汇合起来；第二项要求把两次购买分开。前一项失效，响应丢失后的重试可能再次扣款；后一项失效，用户第二次购买可能被合并。也可以把同身份的请求看作一组：组内请求共同恢复一项操作，组与组之间允许各自产生业务结果。数学上，把对象按这样的相同关系划组，称为建立等价类；在这里，它落实为一条操作记录与引用它的多次尝试。

\begin{figure}[htbp]
\centering
\begin{tikzpicture}[x=1cm,y=1cm]
\draw[draw=Ink!45,fill=Paper,rounded corners=7pt] (0,-2.3) rectangle (5.4,.2);
\draw[draw=Ink!45,fill=Paper,rounded corners=7pt] (7.1,-2.3) rectangle (12.5,.2);
\node[booknote] at (2.7,-.2) {一次购买 $I_1$，身份 $S_1$};
\node[booknote] at (9.8,-.2) {另一次购买 $I_2$，身份 $S_2$};
\foreach \x/\t in {1/R1,2.7/R2,4.4/R3,8.6/R4,11/R5}{
  \node[draw=Ink!55,circle,fill=white,minimum size=9mm,font=\small\sffamily] at (\x,-1.2) {\t};
}
\node[booknote,text width=5.4cm] at (2.7,-3.1)
 {组内身份变化\\同一购买被拆开\\R2 误带新键，可能重复扣款};
\node[booknote,text width=5.4cm] at (9.8,-3.1)
 {两组误用同一身份\\不同购买被合并\\R4 误带旧键，第二份未成交};
\end{tikzpicture}
\caption{稳定性和分离性分别约束组内与组间的身份关系。}
\label{fig:ch03-identity-classes}
\end{figure}


这两项要求说明了稳定身份所承担的责任，也限定了它的能力。客户端把同一次授权分成 K1 和 K2，服务可能合法地受理两个身份；客户端把两次相同商品购买都写成 K1，服务则只能沿用一项记录。服务端的去重规则需要调用方正确分配身份才能对应到用户的授权次数。以后讨论一次操作至多生效一次时，都要保留这项前提。

K1 的生成与保存顺序同样重要。假如客户端先发请求，等收到 P1 后才保存恢复信息，响应丢失时，客户端再次落入第二章的处境。发送前应当先生成 K1，保存它所属的账户、动作、已确认输入和待发送状态。这里的动作指服务要完成的业务种类，例如创建购买；待发送状态说明本地已经登记操作、尚未确认发送结果。客户端确认记录可以恢复后，再发送第一次请求。发生超时，把状态改为结果待确认，并保留 K1。只有操作明确结束或按约定转入人工处理后，才清理恢复入口。

本地持久化完成后、第一次发送前崩溃，恢复程序会找到 K1，再用它发送，请求可能第一次到达服务。第一次发送后崩溃，恢复程序仍会找到 K1，请求可能已完成，也可能尚未执行。服务需要以相同规则接纳这两种到达。客户端的待确认记录，只说明还需要取得远端结果，不能单独证明远端是否提交。

\begin{figure}[htbp]
\centering
\begin{tikzpicture}[x=1cm,y=1cm]
\draw[bookarrow] (0,0)--(13,0);
\foreach \x/\t in {1/生成 K1,4/持久保存操作记录,7.5/发送 R1,11/接收业务结果}{
 \fill[Ink] (\x,0) circle (2pt);
 \node[booknote,text width=3cm,anchor=south] at (\x,.22) {\t};
}
\node[bookresult,text width=3.2cm] (c1) at (5,-1.4) {此处崩溃\\可能尚未发送};
\node[bookresult,text width=3.2cm] (c2) at (9.1,-1.4) {此处崩溃\\远端结果可能未知};
\draw[bookalert] (5,0)--(c1);
\draw[bookalert] (9.1,0)--(c2);
\node[booknode,text width=8.2cm] (r) at (7,-3) {重启后读取原账户、动作、K1 与已确认输入\\用 K1 继续，服务查询或完成同一操作};
\draw[bookarrow] (c1)--(r);
\draw[bookarrow] (c2)--(r);
\end{tikzpicture}
\caption{先保存身份再发送，两处客户端崩溃都能恢复到 K1。}
\label{fig:ch03-client-save}
\end{figure}


\subsection*{页面、进程与设备保存了不同的信息}

只放在内存中的 K1 可以覆盖同一进程里的短暂重试。页面重载、应用被关闭或操作系统回收进程，都可能让它消失。客户端若在恢复后重新生成 K2，服务会收到一项新购买。假设随机数生成器工作得完全正确，也仍然可能产生第二次扣款，因为旧意图已经失去旧身份。

浏览器的本地持久存储可以让同一浏览器恢复 K1。记录应与已登录账户关联，读取时重新核对账户，退出账户后应限制旧记录展示。单独保存一个键还不够：如果页面恢复后用 K1 配上当前购物车，而购物车已经换了商品，服务会收到同键不同输入。恢复记录还应保留确认时的输入、输入解释规则的版本以及必要的报价依据。

多个标签页会引入新的分配问题。两个标签页分别打开同一商品，并各自创建 K1 和 K2，服务会把它们作为两项购买处理。如果产品希望这些页面共同恢复一个待确认操作，就要共享那条本地操作记录，并协调它的创建与选择。但是共享当前商品 ID 仍然无法区分两项合法购买。例如，Alice 上午买一份 \code{paper-001}，下午再买一份，两次确认都对应同一商品 ID；若两个页面只按商品查找本地记录，下午的购买就可能误用上午的 K1。页面需要选择具体的待确认操作，商品 ID 只能帮助展示该操作买了什么。

同时，另一台设备通常读不到第一台设备的本地存储。若产品要求跨设备继续同一购买，可以在服务端保存可由账户查询的购买意图记录，或提供已认证的待处理购买列表，让第二台设备显式选择继续哪一项。设备 B 自己生成的新键只说明它提出了一项新操作，即使它发送的账户与商品和设备 A 完全相同，服务也缺少两者属于同一意图的依据。

设备和页面的恢复都依赖那条操作记录，后台自动重试也应沿用它。假设客户端采用自动重试队列，队列每次取出一条记录，从中读取原键、原商品和报价，刷新访问凭据，再构造新的通信尝试。访问凭据和请求 ID 可以更新，已经确认的业务输入与操作身份继续取自原记录。

账户切换时，还需要暂停这个账户的队列。假设 Alice 的待确认记录属于账户 1，客户端随后登录账户 2；若重试时只读取当前登录凭据，原输入就会被带进账户 2 的识别范围。在账户 2 也有购买权限的情况下，服务甚至可能把它当作账户 2 的新购买处理。队列应核对记录所属账户，重新认证到原账户后再恢复，不能把待确认操作自动迁到新账户。

\begin{figure}[htbp]
\centering
\begin{tikzpicture}[x=1cm,y=1cm]
\node[booknode,text width=4cm] (q) at (2,0) {待确认操作\\账户 1、K1、Q7\\原输入保持不变};
\node[booknode,text width=3.5cm] (auth) at (7,0) {恢复时重新认证\\核对记录所属账户};
\node[booknode,text width=3.4cm] (same) at (11.8,1) {仍为账户 1\\继续 K1};
\node[bookresult,text width=3.4cm] (other) at (11.8,-1) {当前是账户 2\\暂停账户 1 的队列};
\draw[bookarrow] (q)--(auth);
\draw[bookarrow] (auth)--(same);
\draw[bookalert] (auth)--(other);
\node[booknote,text width=12cm] at (6.7,-2.2)
 {可更新访问凭据和本次请求 ID\\账户、身份和已确认输入继续取自原操作};
\end{tikzpicture}
\caption{账户切换需要拦住后台队列，防止把旧购买带入另一账户。}
\label{fig:ch03-queue-account}
\end{figure}


\begin{table}[htbp]
\centering\small
\begin{tabularx}{\linewidth}{@{}p{2.5cm}XX@{}}
\toprule
变化 & 已保存信息的条件 & 恢复动作\\
\midrule
网络超时 & 当前操作仍持有 K1 和输入 & 查询或重发 K1\\
页面重载 & 同账户的本地持久记录仍在 & 从记录恢复 K1\\
客户端重启 & 恢复记录已写入持久介质 & 重新认证后继续 K1\\
多标签页 & 页面共同选择同一操作记录 & 共享 K1 并协调状态\\
更换设备 & 有服务端意图记录或显式交接 & 选择原操作后继续\\
本地记录遗失 & 仅剩账户和商品 & 先核对购买事实，再决定新购买\\
\bottomrule
\end{tabularx}
\caption{不同恢复入口需要保留的身份依据。}
\end{table}

本地数据全部遗失时，客户端查询历史看到一笔 \code{paper-001}，还需要核对它与这次操作的关系：历史中可能原本就有同商品购买。恢复需要更多证据，例如待确认操作记录、明确的成交时间与订单信息，必要时让用户确认。此时如果自动创建 K2 并发送，就会在尚未查清旧购买时增加一项新的扣款授权。

\subsection*{由谁分配身份}

前面一直假定客户端已经拿到 K1。实际设计还要决定谁产生这个值、怎样让不同操作取得不同值，以及分配结果丢失后怎样恢复。身份可以在客户端生成，也可以先向服务申请。两种办法都需要把已选定的值保存下来，否则分配完成后的第一次响应丢失，仍会使下一次尝试丢失原身份。

一种办法由客户端生成足够大的随机键。它无需为分配身份多走一次网络往返，离线时也能准备待发送操作。服务仍要限定键的长度与字符集，核对账户、动作所属的识别范围和参数。独立的新操作可以分别生成键，用户选择重试时则主动沿用原键。

随机生成有一个需要估计的问题：两个独立操作可能偶然取到相同值。键的空间越大，这种相撞越少。一个二进制位可以取 0 或 1，\(b\) 位空间共有 \(2^b\) 种取值。随机熵在这里表示键真正具有多少位独立、不可预测的随机选择。假设每个键独立且均匀地取自这个空间，累计生成 \(n\) 个键；当 \(n\ll 2^{b/2}\) 时，至少出现一次碰撞的概率近似为
\[
 p_{\mathrm{collision}}\approx\frac{n(n-1)}{2^{b+1}}.
\]
这个近似可以从成对比较理解：\(n\) 个键有 \(n(n-1)/2\) 对，每对相同的概率是 \(1/2^b\)。在总碰撞概率很小时，把各对概率相加可以得到上述近似；概率增大后，多个碰撞事件重叠，简单相加就不再准确。若真正具有 128 位随机熵，生成十亿个键，近似概率为 \(1.47\times10^{-21}\)。

这个估计依赖随机源、独立性和实际熵；把时间戳或截短字符串当作 128 位随机数，计算前提就已经改变。第一章使用的 UUID 采用 128 位格式，其中部分位用于标记版本和类别；这些格式位不属于随机熵，计算时应采用所选 UUID 版本的实际随机位数。

概率估计处理的是两个独立新操作偶然取到同一值。客户端把旧 K1 复制给另一项购买，属于确定性的误复用；进程重启后为旧操作生成 K2，则属于身份遗失。增加键长对这两种错误都没有帮助。随机键方案给出的分离保证带有明确的概率前提，稳定重试还依赖恢复记录和客户端操作规则。

还有一种常见生成规则，是以账户编号、商品编号和当天日期拼成键。它把同一天购买同商品的所有操作合并在一起。若再加入秒级时间，用户同一秒内的两次合法购买仍会相撞，旧操作跨秒重试又可能换键。身份生成必须同时满足原操作反复使用相同值、新操作取得不同值两项要求。时间和业务字段可以帮助诊断，还需要明确的新操作入口、保存规则和服务端输入冲突检查。

\begin{figure}[htbp]
\centering
\begin{tikzpicture}[x=1cm,y=1cm]
\node[booknote] at (2,1) {偶然碰撞};
\node[booknote] at (6.6,1) {旧键误复用};
\node[booknote] at (11.2,1) {原身份遗失};
\node[booknode,text width=3.5cm] (a) at (2,0) {新操作 A、B\\分别独立生成键};
\node[booknode,text width=3.5cm] (b) at (6.6,0) {新操作 B\\复制 A 的 K1};
\node[booknode,text width=3.5cm] (c) at (11.2,0) {恢复操作 A\\重新生成 K2};
\node[bookresult,text width=3.5cm] (ae) at (2,-1.7) {偶然得到相同值\\概率依赖键空间};
\node[bookresult,text width=3.5cm] (be) at (6.6,-1.7) {两项购买变成同身份\\更长的 K1 仍被复制};
\node[bookresult,text width=3.5cm] (ce) at (11.2,-1.7)
 {同一购买，两个身份\\随机数生成正确\\仍可能再次扣款};
\draw[bookarrow] (a)--(ae);
\draw[bookalert] (b)--(be);
\draw[bookalert] (c)--(ce);
\end{tikzpicture}
\caption{三个问题依赖不同的规则。增加随机位数主要降低独立生成时的偶然碰撞。}
\label{fig:ch03-key-errors}
\end{figure}


另一种办法由服务预先分配购买意图 token。token 在这里是一段引用服务端意图记录的令牌，客户端随后把它带回来，服务便能找到事先登记的购买准备信息；它是否同时具有授权作用，需要另行规定。客户端先申请一条尚未扣款的意图记录，得到 T1 并保存，随后携带 T1 提交购买。服务可以在记录中绑定账户、动作、报价或截止时间，也可以让多个设备在认证后查询这条意图。代价是多了一次交互，以及意图记录的容量和清理问题。

申请 T1 的响应本身也可能丢失。此时尚未发送购买提交，单纯重复分配 token 不会直接扣款，但可能留下多条空意图。服务可以让分配接口接受一个已持久保存的分配请求键，重复申请返回同一意图；也可以允许空意图自然过期，让客户端在拿到并保存一个 token 后才进入提交阶段。后一办法保住购买身份的前提，是从未拿到并保存 token 的意图不会自动执行购买。

如果 token 分配接口同时扣款，分配响应丢失便重现原问题。把生成编号移到服务端，仍须解释客户端怎样重新找到它。本章采用客户端生成并持久保存 K 的方案继续讨论；有限保存时如何拒绝旧请求，还需要由服务提供可核验的期限依据。
\FloatBarrier
